% ============================================================================
% 模型检验 —— 独立成大章，不是各问里带一句。
% 在同题的人工基线里，这一章是全文第二大的章节；我们历史上的交付物恰恰是
% 把它压缩成几句话，这是页数与 paper 维双输的主因。
%
% 合同（每一小节都要有数字，不能只有论述）：
%   1. 双路互证：同一结论用两条独立路径算（增量 vs 直算、解析 vs 仿真、
%      两套独立实现），给一致性证据表。这一节最容易抓到自己的 bug。
%   2. 与公开/可核事实对表：题面给定的可验证事实、附件能复现的量。
%   3. 参数灵敏度：逐参数扰动 -> 结论是否翻转，给数字不给形容词。
%   4. 样本量与收敛：证书随样本量的变化，说明为什么当前样本量够
%      （相邻档概率差与置信半宽的比较）。
%   5. 稳健性 / 对照臂：换口径、换随机种子、换实现的对照结果。
%   6. 适用边界与失效域：模型在什么条件下不成立，如实写。
% 【v3.2 E/V/A 合同】另须把三路阈值交叉、Balberg/外部理论量级、解析算例与暴力对照、
% 误差传导、多种子临界扫描、有限尺寸/胞元形状、阈值灵敏度系数、全前沿价格灵敏度、
% 上下界夹逼偏差方向、单位步长整数前沿、病态切换和细长胞元反推填入证据矩阵。
% 篇幅参考：5-8 页。
% ============================================================================
\section{模型的检验}

\subsection{双路互证}

\subsection{与可核事实对表}

\subsection{参数灵敏度分析}

\subsection{样本量与收敛性}

\subsection{稳健性与对照实验}

\subsection{适用边界与失效域}

% ---------------------------------------------------------------------------
% 表骨架：表驱动能逼着把数字填出来，比纯论述更难糊弄过去。
% ---------------------------------------------------------------------------

\subsection{误差分析}
\begin{table}[htbp]
  \centering
  \caption{误差来源与量级}
  \begin{tabularx}{\linewidth}{lXll}
    \toprule
    误差来源 & 处理方式（夹逼／对照／声明理由） & 量级 & 对最终答案的影响 \\
    \midrule
     &  &  &  \\
    \bottomrule
  \end{tabularx}
\end{table}
% 【硬条件】最后一列不能空：中间量误差小 ≠ 结论误差小。
% 必须用配对实验把误差传导到最终答案量上，并论证实验分辨力小于所声称的效应量。

\subsection{灵敏度汇总}
\begin{table}[htbp]
  \centering
  \caption{参数灵敏度}
  \begin{tabularx}{\linewidth}{llXl}
    \toprule
    参数 & 扰动幅度 & 输出变化 & 结论是否翻转 \\
    \midrule
     &  &  &  \\
    \bottomrule
  \end{tabularx}
\end{table}

\subsection{E/V/A 证据链汇总}
% 每行必须来自结果/审查/EVA-检查表.md，不能只写「已检查」。
\begin{table}[htbp]
  \centering
  \caption{E/V/A 证据链汇总}
  \begin{tabularx}{\linewidth}{lXlXl}
    \toprule
    维度 & 检查项 & artifact / 命令与退出码 & 数值结果 & 判定 \\
    \midrule
    E & 三路独立估计路径 &  &  &  \\
    E & 外部理论量级（Balberg 或替代） &  &  &  \\
    E/V & 误差传导与临界多 seed &  &  &  \\
    V & 有限尺寸/胞元形状与全前沿价格 &  &  &  \\
    A & 夹逼方向、整数前沿、病态切换、细长胞元 &  &  &  \\
    \bottomrule
  \end{tabularx}
\end{table}
